\section{Von Kuonr\^{a}ts Gefangenschaft und seiner Rückkehr}

\mylettrine{B}{ald} ritt der \gls{Knecht} mit dem Roß \gls{Veltloufer} zurück nach \gls{Hardegg}. Wie mir berichtet ward, vernahm er dort die Kunde vom Ableben des Grafen; kaum war die Nachricht an den Hof gedrungen, so ward der \gls{Knecht} sogleich ausgesandt nach Ottenstein, auf daß dem Vater des Gefangenen Kunde und Beistand zugeleitet werde. Er trug die Botschaft treu und ohne Schmeichelei, denn sein Sinn galt dem Gehorsam, nicht der Sprache.

\gls{Hadmar von Steinare} sandte unverzüglich Boten an den neuen Herzog \gls{Friedrich II von Oesterreich} mit der Bitte um Hilfe, seinen Sohn frei zu bekommen. Mancherlei Umstände, so heißt es, waren dem Herrn Hadmar weder unbekannt noch gering; doch die Welt hat ihre eigenen Drängnisse. Der Herzog war indessen mit einem Aufruhr der \gls{Kueringer} beschäftigt und sammelte Truppen, so daß er die Bitte nicht augenblicklich erhören konnte. Weil Hadmar das Lösegeld, das die Ratsherren von Znaim verlangten, nicht aufzubringen vermochte, blieb der junge Ritter erst in städtischer Verwahrung.

Ein Jahr und ein Tag vergingen, und die Wirren der Nachfolge und des Fehderechts schufen neue Gründe zum Streite. So ward der Aufruhr der \gls{Kueringer} nach mancherlei Gefechten niedergeschlagen, die Feste \gls{Aggstein} fiel, und der Herzog nutzte das gesammelte Heer und den Kriegsgrund, um in das Land Mähren vorzustoßen, sehr zum Missfallen des Königs Wenzel. Manch einer meint, er habe nicht des Kuonr\^{a}ts wegen gehandelt; manch einer schreibt ihm ehrlicher Hoffart zu: Beutegut und die Demonstration seiner Macht. Wie mir berichtet ward, belagerte man auch Znaim; der Rat der Stadt entrichtete eine Summe an Geld und gab Gefangene frei, unter ihnen war auch der junge \gls{Kuonrat von Steinare}.

\gls{Kuonrat von Steinare}s Zeit im Kerker hat sein Gemüt nicht unberührt gelassen. War die Liebe zur \gls{Freiheit} ihm schon zuvor ein inneres Bedürfnis, so ward sie nun zur unverrückbaren Lehre seines Lebens. Die Enge der Zelle, die kalten Speisen und die Demut des Wartens schärften seine Abneigung gegen jede Herrschaft, die ihn bände; zugleich nährte die Ohnmacht vor dem Lösegeld die Habgier, denn Geld schien ihm nunmehr das einzig wirksame Mittel, Ketten zu lösen.

Ob der Herzog in Wahrheit mehr nach Gewinn als nach Gerechtigkeit trachtete, darum streiten die Berichterstatter; manch einer zweifelt daran, doch der Herr allein weiß, welche Gründe im Herzen eines Herrn walten. Für den jungen Kuonr\^{a}t aber ward das Gefängnis zur Schule, und die Sehnsucht nach der freien Fahrt, nach dem Röhren des Rosses und nach dem Gelächter auf Tavernenböden wurde zur harten, unbeugsamen Pflicht: er gelobte sich die \gls{Freiheit} mit jeder Faser seines Willens.

So kam es, wie es kommen mußte: die Rückkehr war kein Triumph der Milde, sondern ein Neuanfang, in dem ein zäher, misstrauischer Sinn reifte. Wie mir berichtet ward, kehrte er heim nicht als ein gebrochener Knabe, sondern als ein Mann, dessen Freiheitswille nunmehr unerschütterlich war. Und manch einer fragte sich, ob das Jahr der Unfreiheit nicht mehr zur Formung seines Charakters beitrug als alle Lehrer zuvor.
